\documentclass[12pt,aspectratio=169,ignorenonframetext]{ltxbeamer}

\usepackage[spanish,es-nodecimaldot,es-tabla]{babel}
\usepackage{hologo}

\title{Grandes Documentos y Programación}
\subtitle{Curso de {\lmr\LaTeX}}
\author{Joar Buitrago\texorpdfstring{\thanks{jebuitragoc@unal.edu.co}}{ <jebuitragoc@unal.edu.co>}}
\date{}

\addbibresource{bibliography.bib}



\begin{document}
\begin{frame}
\titlepage
\end{frame}





\section{Documentos Grandes}

\begin{frame}
\frametitle{El problema de un documento gigante}

Al trabajar en documentos grandes, como por ejemplos tesis o libros, solemos tener una gran cantidad de problemas para llevar un orden dentro del documento.

{\lmr\LaTeX} incluye dos comandos para hacer el manejo de documentos grandes mucho más ameno.
\end{frame}





\begin{frame}[fragile]
\frametitle{El problema de un documento gigante}

\begin{itemize}
    \item \mintinline{latex}{\include}
    \item \mintinline{latex}{\input}
\end{itemize}
\end{frame}





\begin{frame}[fragile]
\frametitle{El problema de un documento gigante}

\begin{minted}[escapeinside=||]{latex}
    \include{|\textit{archivo}|}
\end{minted}

\pause

El comando \mintinline{latex}{\include} hace 3 cosas.

\begin{enumerate}
    \item Realiza una llamada al comando \mintinline{latex}{\clearpage}.
    \item Inserta el código contenido dentro \texttt{\textit{archivo}}.
    \item Realiza una llamada al comando \mintinline{latex}{\clearpage}.
\end{enumerate}
\end{frame}





\begin{frame}[fragile]
\frametitle{El problema de un documento gigante}

El comando \mintinline{latex}{\include} puede ser controlado y manipulado exigiendo que incluya solamente algunos de los elementos escritos.

Esto se logra con el comando \mintinline{latex}{\includeonly}, el cual debe ir en el preámbulo.

\pause

\begin{minted}{latex}
\includeonly{preamble,chapters/chap1}

\include{preamble}
\include{chapters/chap1}
\include{chapters/chap2}
\include{attachments}
\end{minted}

{\lmr\LaTeX} solamente incluirá los archivos \texttt{preamble} y \texttt{chapters/chap1}.
\end{frame}





\begin{frame}[fragile]
\frametitle{El problema de un documento gigante}

\begin{minted}[escapeinside=||]{latex}
    \input{|\textit{archivo}|}
\end{minted}

\pause

El comando \mintinline{latex}{\input} hace una sola cosa.

\begin{enumerate}
    \item Inserta el código contenido dentro \texttt{\textit{archivo}}.
\end{enumerate}
\end{frame}





\begin{frame}[fragile]
\frametitle{El problema de un documento gigante}

Podemos controlar hasta qué lugar del archivo {\lmr\LaTeX} debe incluir, colocando dentro del archivo el comando \mintinline{latex}{\endinput}

\pause

\begin{minted}[breaklines]{latex}
\chapter{Preámbulo}

En este capitulo exploraremos diversos elementos...

\endinput

Esto ya no se incluirá en el documento final.
\end{minted}
\end{frame}





\section{\lmr{\LaTeX} como un lenguaje de programación}

\begin{frame}
\frametitle{{\lmr\LaTeX} como un lenguaje de programación}

{\lmr\TeX} fue diseñado para ser un lenguaje de marcado.
Pero debido a diversas limitantes y complicaciones técnicas que tuvo Donald Knuth durante la creación de este sistema no tuvo más remedio que crear un lenguaje Turing completo.

Esto quiere decir, que {\lmr\TeX} (y por ende {\lmr\LaTeX}) es un lenguaje de programación (no muy eficiente).
\end{frame}





\subsection{{\lmr\TeX}}

\begin{frame}[fragile]
\frametitle{Los comandos \mintinline{latex}{\def} y \mintinline{latex}{\let}}

Estas primitivas de {\lmr\TeX} son bastante útiles para crear comandos sencillos personalizados.
\end{frame}





\begin{frame}[fragile]
\frametitle{Los comandos \mintinline{latex}{\def} y \mintinline{latex}{\let}}

El comando \mintinline{latex}{\def} define un nuevo comando.

\pause

\begin{columns}
    \begin{column}{0.5\textwidth}
        \begin{minted}{latex}
\def\gnuplot{gnuplot}
\def\test{\gnuplot}

\test

\def\gnuplot{not gnuplot}

\test
        \end{minted}
    \end{column}
    {\color{gray} \vline}
    \begin{column}{0.5\linewidth}
        \lmr
        \def\gnuplot{gnuplot}
        \def\test{\gnuplot}

        \test

        \def\gnuplot{not gnuplot}

        \test
    \end{column}
\end{columns}
\end{frame}





\begin{frame}[fragile]
\frametitle{Los comandos \mintinline{latex}{\def} y \mintinline{latex}{\let}}

El comando \mintinline{latex}{\let} asigna una definición.

\pause

\begin{columns}
    \begin{column}{0.5\textwidth}
        \begin{minted}{latex}
\def\gnuplot{gnuplot}
\let\test=\gnuplot

\test

\def\gnuplot{not gnuplot}

\test
        \end{minted}
    \end{column}
    {\color{gray} \vline}
    \begin{column}{0.5\linewidth}
        \lmr
        \def\gnuplot{gnuplot}
        \let\test=\gnuplot

        \test

        \def\gnuplot{not gnuplot}

        \test
    \end{column}
\end{columns}
\end{frame}





\begin{frame}[fragile]
\frametitle{Los comandos \mintinline{latex}{\def} y \mintinline{latex}{\let}}

\begin{alertblock}{IMPORTANTE}
    Estas primitivas de {\lmr\TeX} no tienen en cuenta si ya existía definición previa, por lo que pueden generar errores si se redefinen comandos propios de {\lmr\TeX} o {\lmr\LaTeX}.
\end{alertblock}
\end{frame}





\begin{frame}[fragile]
\frametitle{{\lmr\LaTeXe}}
\framesubtitle{Comando}

\begin{itemize}
    \item \mintinline{latex}{\newcommand}
    \item \mintinline{latex}{\renewcommand}
    \item \mintinline{latex}{\providecommand}
\end{itemize}
\end{frame}





\subsection{{\lmr\LaTeXe}}

\begin{frame}[fragile]
\frametitle{{\lmr\LaTeXe}}
\framesubtitle{Comando}

\begin{minted}[escapeinside=||]{latex}
\newcommand
    *
    {\|\textit{comando}|}
    [|\textit{núm. argumentos}|]
    [|\textit{valor por defecto}|]
    {|\textit{definición}|}
\end{minted}
\end{frame}





\begin{frame}[fragile]
\frametitle{{\lmr\LaTeXe}}
\framesubtitle{Comando}

\begin{itemize}
    \item \mintinline{latex}{\newcommand}: Genera un error si el comando ya existe.
    \item \mintinline{latex}{\renewcommand}: Arroja un error si el comando no existía.
    \item \mintinline{latex}{\providecommand}: No redefine el comando si este ya existe.
\end{itemize}
\end{frame}





\begin{frame}[fragile]
\frametitle{{\lmr\LaTeXe}}
\framesubtitle{Comando}

\begin{columns}
    \begin{column}{0.5\textwidth}
        \begin{minted}{latex}
\newcommand
    {\test}
    [1]
    [Mundo]
    {Hola #1}

\test

\test[German]
        \end{minted}
    \end{column}
    {\color{gray} \vline}
    \begin{column}{0.5\linewidth}
        \lmr
        \newcommand{\test}[1][Mundo]{Hola #1}

        \test

        \test[German]
    \end{column}
\end{columns}
\end{frame}





\begin{frame}[fragile]
\frametitle{{\lmr\LaTeXe}}
\framesubtitle{Entorno}

\begin{itemize}
    \item \mintinline{latex}{\newenvironment}\only<2->{: Genera un error si el entorno ya existe.}
    \item \mintinline{latex}{\renewenvironment}\only<2->{: Arroja un error si el entorno no existía.}
    \item \mintinline{latex}{\provideenvironment}\only<2->{: No redefine el entorno si este ya existe.}
\end{itemize}
\end{frame}





\begin{frame}[fragile]
\frametitle{{\lmr\LaTeXe}}
\framesubtitle{Entorno}

\begin{minted}[escapeinside=||]{latex}
\newenvironment
    *
    {|\textit{entorno}|}
    [|\textit{núm. argumentos}|]
    [|\textit{valor por defecto}|]
    {|\textit{principio}|}
    {|\textit{final}|}
\end{minted}
\end{frame}





\begin{frame}[fragile]
\frametitle{{\lmr\LaTeXe}}
\framesubtitle{Entorno}

\begin{columns}
    \begin{column}{0.5\textwidth}
        \begin{minted}{latex}
\newenvironment
    {test}
    {Hola}
    {y Adios}

\begin{test}
    Mundo
\end{test}
        \end{minted}
    \end{column}
    {\color{gray} \vline}
    \begin{column}{0.5\linewidth}
        \lmr
        \newenvironment
            {test}
            {Hola}
            {y Adios}

        \begin{test}
            Mundo
        \end{test}
    \end{column}
\end{columns}
\end{frame}





\begin{frame}[fragile]
\frametitle{{\lmr\LaTeXe}}
\framesubtitle{Contador}

\begin{itemize}
    \item \mintinline{latex}{\newcounter}
    \item \mintinline{latex}{\stepcounter}
    \item \mintinline{latex}{\setcounter}
    \item \mintinline{latex}{\addtocounter}
\end{itemize}
\end{frame}





\begin{frame}[fragile]
\frametitle{{\lmr\LaTeXe}}
\framesubtitle{Contador}

\begin{columns}
    \begin{column}{0.5\textwidth}
        \begin{minted}{latex}
\newcounter{mycounter}
\arabic{mycounter}\par
\setcounter{mycounter}{16}
\alph{mycounter}\par
\addtocounter{mycounter}{5}
\roman{mycounter}\par
\addtocounter{mycounter}{-3}
\Alph{mycounter}\par
\addtocounter{mycounter}
    {\value{mycounter}}
\Roman{mycounter}\par
\stepcounter{mycounter}
\arabic{mycounter}\par
        \end{minted}
    \end{column}
    {\color{gray} \vline}
    \begin{column}{0.5\linewidth}
        \lmr
        \newcounter{mycounter}
        \arabic{mycounter}\par
        \setcounter{mycounter}{16}
        \alph{mycounter}\par
        \addtocounter{mycounter}{5}
        \roman{mycounter}\par
        \addtocounter{mycounter}{-3}
        \Alph{mycounter}\par
        \addtocounter{mycounter}{\value{mycounter}}
        \Roman{mycounter}\par
        \stepcounter{mycounter}
        \arabic{mycounter}\par
    \end{column}
\end{columns}
\end{frame}





\begin{frame}[fragile]
\frametitle{{\lmr\LaTeXe}}
\framesubtitle{Contador}

\begin{columns}
    \begin{column}{0.5\textwidth}
        \begin{minted}{latex}
\newcounter{parentcounter}
\newcounter{childcounter}
    [parentcounter]
\stepcounter{parentcounter}
\stepcounter{childcounter}
\arabic{parentcounter}-
    \arabic{childcounter}\par
\stepcounter{childcounter}
\stepcounter{childcounter}
\arabic{parentcounter}-
    \arabic{childcounter}\par
\stepcounter{parentcounter}
\stepcounter{childcounter}
\arabic{parentcounter}-
    \arabic{childcounter}\par
        \end{minted}
    \end{column}
    {\color{gray} \vline}
    \begin{column}{0.5\linewidth}%
        \lmr%
        \newcounter{parentcounter}%
        \newcounter{childcounter}[parentcounter]%
        \stepcounter{parentcounter}%
        \stepcounter{childcounter}%
        \arabic{parentcounter}-\arabic{childcounter}\par
        \stepcounter{childcounter}%
        \stepcounter{childcounter}%
        \arabic{parentcounter}-\arabic{childcounter}\par
        \stepcounter{parentcounter}%
        \stepcounter{childcounter}%
        \arabic{parentcounter}-\arabic{childcounter}\par
    \end{column}
\end{columns}
\end{frame}





\begin{frame}[fragile]
\frametitle{{\lmr\LaTeXe}}
\framesubtitle{Longitud}

\begin{itemize}
    \item \mintinline{latex}{\newlength}
    \item \mintinline{latex}{\setlength}
    \item \mintinline{latex}{\addtolength}
    \item \mintinline{latex}{\settodepth}
    \item \mintinline{latex}{\settoheight}
    \item \mintinline{latex}{\settowidth}
\end{itemize}
\end{frame}





\begin{frame}[fragile]
\frametitle{{\lmr\LaTeXe}}
\framesubtitle{Longitud}

\begin{columns}
    \begin{column}{0.5\textwidth}
        \begin{minted}{latex}
\newlength{\mylength}
\the\mylength\par
\setlength{\mylength}{1em}
\the\mylength\par
\addtolength{\mylength}{10ex}
\the\mylength\par
\addtolength{\mylength}{-5mm}
\the\mylength\par
        \end{minted}
    \end{column}
    {\color{gray} \vline}
    \begin{column}{0.5\linewidth}
        \lmr%
        \newlength{\mylength}%
        \the\mylength\par
        \setlength{\mylength}{1em}%
        \the\mylength\par
        \addtolength{\mylength}{10ex}%
        \the\mylength\par
        \addtolength{\mylength}{-5mm}%
        \the\mylength\par
    \end{column}
\end{columns}
\end{frame}





\begin{frame}[fragile]
\frametitle{{\lmr\LaTeXe}}
\framesubtitle{Longitud}

\begingroup
\centering
\begin{tikzpicture}
    \node [inner sep=0pt] (O) {\fontsize{56}{60}\lmr Murcielago};
    \draw [|<->|] ([shift={(0,0.5)}]O.north west) -- ([shift={(0,0.5)}]O.north east) node [midway, above] {\mintinline{latex}{\width}};
    \draw [|<->|] ([shift={(-0.5,0)}]O.north west) -- ([shift={(-0.5,0)}]O.south west) node [midway, left] {\mintinline{latex}{\totalheight}};
    \draw [|<->|] ([shift={(0.5,0)}]O.north east) -- ([shift={(0.5,0)}]O.north east |- O.base) node [midway, right] {\mintinline{latex}{\height}};
    \draw [|<->|] ([shift={(0.5,0)}]O.north east |- O.base) -- ([shift={(0.5,0)}]O.south east) node [midway, right] {\mintinline{latex}{\depth}};
\end{tikzpicture}
\par
\endgroup
\end{frame}





\begin{frame}[fragile]
\frametitle{{\lmr\LaTeXe}}
\framesubtitle{Booleanos}

\mintinline{latex}{\newif}
\end{frame}





\begin{frame}[fragile]
\frametitle{{\lmr\LaTeXe}}
\framesubtitle{Longitud}

\begin{columns}
    \begin{column}{0.5\textwidth}
        \begin{minted}{latex}
\newif\ifworks
\workstrue
\ifworks
    Es verdad
\else
    Es falso
\fi
        \end{minted}
    \end{column}
    {\color{gray} \vline}
    \begin{column}{0.5\linewidth}
        \lmr%
        \newif\ifworks
        \workstrue
        \ifworks
            Es verdad
        \else
            Es falso
        \fi
    \end{column}
\end{columns}
\end{frame}





\subsection{{\lmr\hologo{LaTeX3}}}

\begin{frame}[fragile]
\frametitle{{\lmr\hologo{LaTeX3}}}
\framesubtitle{Comando}

\begin{itemize}
    \item \mintinline{latex}{\NewDocumentCommand}: Genera un error si el comando ya existe.
    \item \mintinline{latex}{\RenewDocumentCommand}: Arroja un error si el comando no existía.
    \item \mintinline{latex}{\ProvideDocumentCommand}: No redefine el comando si este ya existe.
    \item \mintinline{latex}{\DeclareDocumentCommand}: Redefine siempre el comando.
\end{itemize}
\end{frame}





\begin{frame}[fragile]
\frametitle{{\lmr\hologo{LaTeX3}}}
\framesubtitle{Comando}

\begin{minted}[escapeinside=||]{latex}
\NewDocumentCommand
    {\|\textit{comando}|}
    {|\textit{especificador de argumentos}|}
    {|\textit{definición}|}
\end{minted}
\end{frame}





\begin{frame}[fragile]
\frametitle{{\lmr\hologo{LaTeX3}}}
\framesubtitle{Comando}

\begin{columns}
    \begin{column}{0.5\textwidth}
        \begin{minted}{latex}
\NewDocumentCommand
    {\salutation}
    { O{Mundo} }
    {Hola #1}

\salutation

\salutation[German]
        \end{minted}
    \end{column}
    {\color{gray} \vline}
    \begin{column}{0.5\linewidth}
        \lmr
        \NewDocumentCommand
            {\salutation}
            { O{Mundo} }
            {Hola #1}

        \salutation

        \salutation[German]
    \end{column}
\end{columns}
\end{frame}





\begin{frame}[fragile]
\frametitle{{\lmr\hologo{LaTeX3}}}
\framesubtitle{Comando}

\begin{columns}
    \begin{column}{0.5\textwidth}
        \begin{minted}{latex}
\NewDocumentCommand
    {\salutation}
    { O{Hola} O{Mundo} }
    {#1 #2}

\salutation

\salutation[Saluton]

\salutation[Saluton][German]
        \end{minted}
    \end{column}
    {\color{gray} \vline}
    \begin{column}{0.5\linewidth}
        \lmr
        \NewDocumentCommand
            {\salutation}
            { O{Hola} O{Mundo} }
            {#1 #2}

        \salutation

        \salutation[Saluton]

        \salutation[Saluton][German]
    \end{column}
\end{columns}
\end{frame}






\begin{frame}[fragile]
\frametitle{{\lmr\hologo{LaTeX3}}}
\framesubtitle{Entorno}

\begin{itemize}
    \item \mintinline{latex}{\NewDocumentEnvironment}: Genera un error si el entorno ya existe.
    \item \mintinline{latex}{\RenewDocumentEnvironment}: Arroja un error si el entorno no existía.
    \item \mintinline{latex}{\ProvideDocumentEnvironment}: No redefine el entorno si este ya existe.
    \item \mintinline{latex}{\DeclareDocumentEnvironment}: Redefine siempre el entorno.
\end{itemize}
\end{frame}





\begin{frame}[fragile]
\frametitle{{\lmr\hologo{LaTeX3}}}
\framesubtitle{Entorno}

\begin{minted}[escapeinside=||]{latex}
\NewDocumentEnvironment
    {|\textit{entorno}|}
    {|\textit{especificador de argumentos}|}
    {|\textit{principio}|}
    {|\textit{final}|}
\end{minted}
\end{frame}





\begin{frame}[fragile]
\frametitle{{\lmr\hologo{LaTeX3}}}
\framesubtitle{Entorno}

\begin{columns}
    \begin{column}{0.5\textwidth}
        \begin{minted}{latex}
\NewDocumentEnvironment
    {welcome}
    {}
    {Hola}
    {y Adios}

\begin{welcome}
    Mundo
\end{welcome}
        \end{minted}
    \end{column}
    {\color{gray} \vline}
    \begin{column}{0.5\linewidth}
        \lmr
        \NewDocumentEnvironment
            {welcome}
            {}
            {Hola}
            {y Adios}

        \begin{welcome}
            Mundo
        \end{welcome}
    \end{column}
\end{columns}
\end{frame}




\begin{frame}
\frametitle{{\lmr\hologo{LaTeX3}}}
\framesubtitle{Especificadores obligatorios}

\setbeamersize{description width={0.4\textwidth}}
\begin{description}
    \item[\texttt{m}] Un argumento \emph{mandatorio}.
    \item[\texttt{r \textit{tokenA} \textit{tokenB}}] Un argumento \emph{requerido}, rodeado por \texttt{\textit{tokenA}} y \texttt{\textit{tokenB}}.
    \item[\texttt{R \textit{tokenA} \textit{tokenB} \{\textit{default}\}}] Un argumento \emph{requerido}, rodeado por \texttt{\textit{tokenA}} y \texttt{\textit{tokenB}}.
    \item[\texttt{v}] Un argumento \emph{verbatim} o literal, no será procesado por {\lmr\LaTeX} como código.
    \item[\texttt{b}] En un entorno, los contenidos del entorno.
    \item[\texttt{c}] En un entorno, los contenidos literales del entorno.
\end{description}
\end{frame}





\begin{frame}
\frametitle{{\lmr\hologo{LaTeX3}}}
\framesubtitle{Especificadores opcionales}

\setbeamersize{description width={0.4\textwidth}}
\begin{description}
    \item[\texttt{o}] Un argumento \emph{opcional}.
    \item[\texttt{d \textit{tokenA} \textit{tokenB}}] Un argumento \emph{opcional}, rodeado por \texttt{\textit{tokenA}} y \texttt{\textit{tokenB}}.
    \item[\texttt{O \{\textit{default}\}}] Un argumento \emph{opcional} con valor por defecto.
    \item[\texttt{D \textit{tokenA} \textit{tokenB} \{\textit{default}\}}] Un argumento opcional, rodeado por \texttt{\textit{tokenA}} y \texttt{\textit{tokenB}}, con valor por defecto.
    \item[\texttt{s}] Un comando con posibilidad de ser estrellado.
    \item[\texttt{t \textit{token}}] Un comando con posibilidad de ser \emph{token}izado.
\end{description}
\end{frame}





\section{Clases y Paquetes}

\begin{frame}
\frametitle{Clases y Paquetes}

Las clases y los paquetes de {\lmr\LaTeX} se definen en archivos \texttt{.cls} y \texttt{.sty}.

\pause

Estos archivos {\lmr\LaTeX} los busca en:

\begin{itemize}
    \item El directorio de instalación de la distribución y subdirectorios.
    \item El directorio actual.
    \item Los directorios definidos en la variable de entorno \texttt{TEXINPUTS}.
\end{itemize}
\end{frame}





\begin{frame}[fragile]
\frametitle{Estructura de una clase}

\begin{minted}{latex}
\NeedsTeXFormat{LaTeX2e}

\ProvidesClass{article}[1806-08-20 Article class by OAN]

Código de la clase
\end{minted}
\end{frame}





\begin{frame}[fragile]
\frametitle{Estructura de un paquete}

\begin{minted}{latex}
\NeedsTeXFormat{LaTeX2e}

\ProvidesPackage{gnuplot}[1806-08-20 gnuplot support]

Código de la clase
\end{minted}
\end{frame}





\begin{frame}[fragile]
\frametitle{Opciones}

\begin{minted}[escapeinside=||]{latex}
\DeclareOption
    {|\textit{opcion}|}
    {|\textit{código}|}
\end{minted}
\end{frame}





\begin{frame}[fragile]
\frametitle{Opciones}

\begin{minted}[escapeinside=||]{latex}
\DeclareOption*
    {|\textit{código}|}
\end{minted}
\end{frame}





\begin{frame}[fragile]
\frametitle{Opciones}

\begin{minted}[escapeinside=||]{latex}
\DeclareOption{color}{\colortrue}

\DeclareOption*{\PassOptionsToClass{\CurrentOption}{article}}
\end{minted}
\end{frame}





\begin{frame}[fragile]
\frametitle{Opciones}

\begin{minted}[escapeinside=||]{latex}
\ProcessOptions\relax
\end{minted}
\end{frame}





\begin{frame}[fragile]
\frametitle{Opciones}

\begin{minted}[escapeinside=||]{latex}
\LoadClass{article}
\end{minted}
\end{frame}





\begin{frame}[fragile]
\frametitle{Opciones}

\begin{minted}[escapeinside=||]{latex}
\RequirePackage{geometry}
\end{minted}
\end{frame}





\begin{frame}[fragile]
\frametitle{{\lmr\hologo{LaTeX3}}}
\framesubtitle{Entorno}

\begin{columns}
    \begin{column}{0.5\linewidth}
        \lmr
        \newif\ifprime \newif\ifunknown %
        \newcount\n \newcount\p \newcount\d \newcount\a %
        \def\primes#1{2,~3%
            \n=#1 \advance\n by-2 %
            \p=5 %
            \loop\ifnum\n>0 \printifprime\advance\p by2 \repeat} %
        \def\printp{, %
            \ifnum\n=1 y\ \fi %
            \number\p \advance\n by -1 }
            \def\printifprime{\testprimality \ifprime\printp\fi} %
        \def\testprimality{{\d=3 \global\primetrue
            \loop\trialdivision \ifunknown\advance\d by2 \repeat}} %
        \def\trialdivision{\a=\p \divide\a by\d
            \ifnum\a>\d \unknowntrue\else\unknownfalse\fi %
            \multiply\a by\d %
            \ifnum\a=\p \global\primefalse\unknownfalse\fi} %

        Los primeros 30 primos son \primes{30}.
    \end{column}
    \pause
    {\color{gray} \vline}
    \begin{column}{0.5\textwidth}
        \tiny
        \begin{minted}[breaklines]{latex}
\newif\ifprime \newif\ifunknown %
\newcount\n \newcount\p \newcount\d \newcount\a %
\def\primes#1{2,~3%
    \n=#1 \advance\n by-2 %
    \p=5 %
    \loop\ifnum\n>0 \printifprime\advance\p by2 \repeat}
\def\printp{, %
    \ifnum\n=1 y\ \fi %
    \number\p \advance\n by -1 }
    \def\printifprime{\testprimality \ifprime\printp\fi}
\def\testprimality{{\d=3 \global\primetrue
    \loop\trialdivision \ifunknown\advance\d by2 \repeat}}
\def\trialdivision{\a=\p \divide\a by\d
    \ifnum\a>\d \unknowntrue\else\unknownfalse\fi
    \multiply\a by\d
    \ifnum\a=\p \global\primefalse\unknownfalse\fi}

Los primeros 30 primos son \primes{30}.
        \end{minted}
    \end{column}
\end{columns}
\end{frame}





\begin{frame}[allowframebreaks]
\frametitle{Referencias}
\nocite{*}
\printbibliography
\end{frame}
\end{document}
